\section{Introduction}

Graph neural networks (GNNs) have become the tool of choice for learning on the structured data that pervades modern applications -- network flows, system call trees, privilege graphs, molecules, and social networks are all most faithfully represented as graphs. Yet the practical performance of a GNN hinges on hyperparameter optimization (HPO), an often overlooked but decisive step: the choice of learning rate, regularization strength, hidden dimensions, or sampling strategy routinely swings model quality by wide margins. Because each candidate configuration requires training a GNN to convergence, HPO is also expensive, and this cost is precisely what motivates methods that can transfer knowledge from previously tuned datasets to new ones.

\subsection{Related Work}\label{sec:related-work}

HPO for graph learning remains a young field, driven by the high cost of GNN training and the complicated interactions between their hyperparameters. A large share of the existing effort falls under Neural Architecture Search (NAS) for GNNs, where GraphNAS \cite{gao_graph_2020} and Auto-GNN \cite{zhou_auto-gnn_2022} use reinforcement learning to search the space of GNN architectures. These methods search each dataset from scratch, however, and neither reuses the results of prior searches nor exploits the structure shared across datasets—the very information a pre-trained surrogate is designed to capture.

Studies that tune GNN hyperparameters directly tend to be confined to a single application domain, such as molecular property prediction \cite{yuan_genetic_2021, yuan_systematic_2021, yuan_which_2021, ebadi_hyperparameter_2025} or mRNA degradation prediction \cite{vodilovska_hyperparameter_2023}. Because they draw on only a handful of datasets from one domain, they offer little basis for learning how tuning outcomes transfer across the diverse range of graph structures encountered in practice—the setting our synthetic pre-training approach is meant to address.

Pretraining on synthetic data is not a novel idea -- it is one of the core ideas behind recent works such as TabPFN \cite{hollmann_accurate_2025}. However, this idea has not been explored in the context of hyperparameter optimization.

\subsection{Contributions}

Our previous work \cite{dedic_benchmarking_2026} studies HPO for general GNNs, benchmarks the most commonly used HPO methods for GraphSAGE \cite{hamilton_inductive_2017} across a diverse set of graph datasets, and introduces a meta-learning approach for hyperparameter transfer. This work extends that by looking at the hyperparameter transfer problem as a pre-training problem, where we learn a surrogate model that can predict the performance of hyperparameter configurations on unseen datasets based on their meta-features. This allows us to warm-start the HPO process on new datasets, potentially reducing the number of costly evaluations needed to find optimal hyperparameters. In \cite{dedic_benchmarking_2026}, this was done on other real-world datasets, limiting its breadth. In this work, we instead use synthetic datasets in the pre-training phase, which allows us to be more flexible in the pre-training and not be limited by the available real-world datasets.

In Section~\ref{sec:related-work}, we reviewed related work on HPO methods, particularly in the context of graph learning. Section~\ref{sec:problem-definition} provides necessary context of HPO more generally. Section~\ref{sec:cross-rf} recalls the Cross-RF HPO method for graph learning from our previous work and describes its limitations. Section~\ref{sec:synthetic-rf} introduces the Synthetic-RF method, which addresses these limitations by pre-training a surrogate model on synthetic datasets. Section~\ref{sec:experimental-evaluation} outlines our methodology, hyperparameter spaces, and evaluated methods and reports on the performance of the Synthetic-RF method, comparing its performance to random search. Finally, Section~\ref{sec:conclusion} summarizes our findings and discusses future research directions.
